Throughout, indices are taken modulo \(n\).

\textbf{Setup.}
Write \(\varepsilon_i := \sigma(t_i)\in\{+1,-1\}\).
Since \(t_{i+1}=t_i+m_ia_i\) and the signs alternate along each chain \(F_i\),
\[
  \varepsilon_{i+1} = (-1)^{m_i}\,\varepsilon_i.
\]
Going once around the full cycle,
\[
  \varepsilon_1 = (-1)^{m_1+\cdots+m_n}\,\varepsilon_1,
\]
so \(\sum_{i=1}^n m_i \equiv 0\pmod{2}\) \textbf{(parity)}.  Summing the edge relations around the cycle gives
\[
  \sum_{i=1}^n m_i a_i = 0 \qquad\textbf{(closure)}.
\]

Under the extremal equality, the lonely-point analysis from Subproblem 1 gives
\[
  \operatorname{supp}(\sigma)=\{x_i=s_i+t_i:1\le i\le n\}
\]
with \(\sigma(x_i)=\varepsilon_i\).  In particular, for every \(t\in T\) and every \(s_j\in S\),
\begin{equation}\label{eq:cancel}
  \sigma(s_j+t)=0 \quad\text{if }s_j+t\notin\{x_1,\dots,x_n\},
  \tag{$*$}
\end{equation}
and \(s_j+t=x_k\) iff \(t=t_k+(s_k-s_j)\) (possible for at most one \(k\)).

\textbf{Remark on both signs.}
Since both signs are present in \(T\) and \(|T|\ge 2\), we have \(m_i\ge 1\) for at least one \(i\) (otherwise every chain \(F_i\) is the single point \(t_i\), and since the chains share endpoints the entire set $T$ would contain only one translate, contradicting both signs being present).

---

\section*{1. Triangle case (\(n=3\))}

The vertex sum gives \(a_1+a_2+a_3=0\), so \(a_3=-(a_1+a_2)\), and \(a_1,a_2\) are linearly independent (three distinct direction classes means no two \(a_i\) are scalar multiples of each other).  The closure condition
\[
  m_1 a_1+m_2 a_2+m_3 a_3=0
  \implies (m_1-m_3)a_1+(m_2-m_3)a_2=0
  \implies m_1=m_2=m_3=:m.
\]
Parity: \(3m\equiv 0\pmod 2\), so \(m\) is even.  Since both signs are present and \(F_i\) is the \emph{only} element of \(T\) if \(m=0\) (giving \(T=\{t_1\}\), one sign), we have \(m\ge 2\).

Because \(m\) is even, \(\varepsilon_2=(-1)^m\varepsilon_1=\varepsilon_1\) and similarly \(\varepsilon_3=\varepsilon_1\).  Write \(\varepsilon:=\varepsilon_1=\varepsilon_2=\varepsilon_3\).

\subsection*{1.1 Net sign from each chain}

Chain \(F_i=\{t_i,t_i+a_i,\dots,t_i+ma_i\}\) has \(m+1\) elements with alternating signs starting at \(\varepsilon\).  Since \(m\) is even, positions \(0,2,\dots,m\) (count: \(m/2+1\)) have sign \(\varepsilon\) and positions \(1,3,\dots,m-1\) (count: \(m/2\)) have sign \(-\varepsilon\).  Net:
\[
  \sum_{t\in F_i}\sigma(t) = \bigl(\tfrac m2+1\bigr)\varepsilon - \tfrac m2\,\varepsilon = \varepsilon.
\]
The endpoints of consecutive chains coincide: \(t_2\in F_1\cap F_2\), \(t_3\in F_2\cap F_3\), \(t_1\in F_3\cap F_1\), each with sign \(\varepsilon\).  By inclusion--exclusion:
\[
  \sum_{t\in F_1\cup F_2\cup F_3}\!\sigma(t)
  = \varepsilon+\varepsilon+\varepsilon - \sigma(t_2)-\sigma(t_3)-\sigma(t_1)
  = 3\varepsilon - 3\varepsilon = 0.
\]

\subsection*{1.2 Staircase forces $T\subseteq F_1\cup F_2\cup F_3$}

Write every element \(t\in T\) in the basis \(\{a_1,a_2\}\) as
\[
  t = t_1 + \alpha(t)\,a_1 + \beta(t)\,a_2.
\]
(This is unique since \(a_1,a_2\) is a basis for \(\mathbb{R}^2\).)

For any \(t_0\in T\), expanding \(\sigma(s_3+t_0)\) using $(\ast)$:
\[
  \sigma(s_3+t_0)
  = \sigma(t_0)\cdot\mathbf{1}_{t_0\in T}
  + \sigma(t_0+a_2)\cdot\mathbf{1}_{t_0+a_2\in T}
  + \sigma(t_0+a_1+a_2)\cdot\mathbf{1}_{t_0+a_1+a_2\in T},
\]
since \(s_3-s_2=a_2\) and \(s_3-s_1=a_1+a_2\), so the three possible values of \(p-s_j\) for \(p=s_3+t_0\) are \(t_0\) (\(j=3\)), \(t_0+a_2\) (\(j=2\)), \(t_0+a_1+a_2\) (\(j=1\)).

We determine when \(\sigma(s_3+t_0)\ne 0\): this happens iff \(s_3+t_0\in\{x_1,x_2,x_3\}\), i.e.\ \(t_0\in\{t_j+(s_j-s_3):j=1,2,3\}\):
\begin{align*}
  t_0 = t_1+s_1-s_3 &= t_1-(a_1+a_2),\quad\beta = -1,\\
  t_0 = t_2+s_2-s_3 &= t_1+ma_1-a_2,\quad\beta = -1,\\
  t_0 = t_3 &= t_1+ma_1+ma_2,\quad\beta = m.
\end{align*}

\medskip\noindent
\textbf{Claim.}  The maximum value of \(\beta(t)\) over all \(t\in T\) is \(m\), achieved only by \(t_3\).

\smallskip\noindent\textit{Proof of Claim.}
Let \(\beta_{\max}=\max_{t\in T}\beta(t)\).  Let \(t^*\in T\) achieve this maximum.  Then \(t^*+a_2\notin T\) and \(t^*+a_1+a_2\notin T\) (both have \(\beta=\beta_{\max}+1>\beta_{\max}\)).  Thus \(\sigma(s_3+t^*)=\sigma(t^*)\ne 0\), so \(t^*\) is one of the three escape points above, giving \(\beta(t^*)\in\{-1,m\}\).  Since \(t_3\in T\) has \(\beta(t_3)=m\), we have \(\beta_{\max}\ge m\).  So \(\beta_{\max}=m\) and \(t^*=t_3\).\hfill\(\square\)

\medskip\noindent
\textbf{Staircase lemma.}  For every \(t_0\in T\) with \(\sigma(s_3+t_0)=0\):
\begin{enumerate}
  \item Exactly one of \(\{t_0+a_2,\,t_0+a_1+a_2\}\) belongs to \(T\).
  \item That element has sign \(-\sigma(t_0)\).
\end{enumerate}

\smallskip\noindent\textit{Proof.}
We have \(\sigma(t_0)+\sigma(t_0+a_2)\cdot\mathbf{1}+\sigma(t_0+a_1+a_2)\cdot\mathbf{1}=0\).  Since each sign is \(\pm 1\) when present: two elements of \(\{\pm 1\}\) summing to \(-\sigma(t_0)\in\{\pm 1\}\) is impossible (their sum would be \(0\) or \(\pm 2\)).  So exactly one is in \(T\), with sign \(-\sigma(t_0)\).\hfill\(\square\)

\medskip\noindent
\textbf{Staircase argument.}  Suppose \(t_0\in T\) with \(\sigma(s_3+t_0)=0\).  Define the staircase sequence \(t_0,t_1,t_2,\dots\) by: \(t_{k+1}\) is the unique element of \(\{t_k+a_2,t_k+a_1+a_2\}\cap T\).  This is well-defined as long as \(\sigma(s_3+t_k)=0\).

At each step, \(\beta(t_{k+1})=\beta(t_k)+1\) (both choices increase \(\beta\) by 1).  By the Claim, \(\beta(t)\le m\) for all \(t\in T\).  So the sequence is finite.

Termination: the sequence ends when \(\sigma(s_3+t_K)\ne 0\) for the first time, meaning \(\beta(t_K)\in\{-1,m\}\).  Since \(\beta(t_k)=\beta(t_0)+k\) and \(\beta(t_0)\ge -1\) (if \(t_0\in T\)), and the sequence is strictly increasing, the only valid termination is \(\beta(t_K)=m\), i.e.\ \(t_K=t_3\).

\medskip\noindent
\textbf{Conclusion: }$T\subseteq F_1\cup F_2\cup F_3$.

Elements of \(F_2=\{t_1+ma_1+ka_2:0\le k\le m\}\) have \(\beta=k\in\{0,\dots,m\}\), \(\alpha=m\).

Elements of \(F_3=\{t_1+ja_1+ja_2:0\le j\le m\}\) (since \(a_3=-a_1-a_2\), so \(t_3+ka_3=t_1+(m-k)(a_1+a_2)\), set \(j=m-k\)) have \(\alpha=\beta=j\).

At any staircase level \(\beta=k\) (for \(0\le k<m\)), the two candidate successors for any element \(({\alpha},k)\in T\) are \((\alpha,k+1)\) and \((\alpha+1,k+1)\).  We claim:  any element \((\alpha,k)\) with \(0<\alpha<m\) and \(\alpha\ne k\) (a ``middle'' element not in \(\cup F_i\)) leads to a contradiction.

If \((\alpha,k)\in T\) (meaning \(t=t_1+\alpha a_1+ka_2\in T\)) with \(\alpha\ne 0,m\) and \(\alpha\ne k\):  by the staircase, its two possible successors are \(({\alpha},k+1)\) and \((\alpha+1,k+1)\).  Continue the staircase from both: eventually it must reach \(t_3=(\alpha=m,\beta=m)\).  At each level, the only elements in \(\cup F_i\) at \(\beta=j\) are \((m,j)\in F_2\) and \((j,j)\in F_3\).

Now suppose the staircase from \(t=(\alpha,k)\) (middle, \(\sigma(s_3+t)=0\)) hits a level \(\beta=j\) where both candidates \((\alpha',j)\) and \((\alpha'+1,j)\) are in \(\cup F_i\subseteq T\).  This happens when \(\alpha'=m-1\) (giving \((m-1,j)\notin \cup F_i\) but \((m,j)\in F_2\)) or \(\alpha'=j-1\) (giving \((j-1,j)\notin \cup F_i\) but \((j,j)\in F_3\)).

The critical case: both candidates at level \(\beta=j+1\) are in \(T\).  Then the staircase equation at \(t'\) (level \(\beta=j\)):
\[
  \sigma(t') + \sigma(t'+a_2) + \sigma(t'+a_1+a_2) = 0
  \implies \sigma(t') = -\sigma(t'+a_2)-\sigma(t'+a_1+a_2).
\]
If both \(t'+a_2\) and \(t'+a_1+a_2\) are in \(T\) with the same sign (forced by them being chain elements with the same \(\beta\)-level sign from \(\cup F_i\)), then \(\sigma(t')=-2\cdot(\pm 1)\notin\{+1,-1\}\).  Contradiction.

(More precisely: elements of \(F_2\) and \(F_3\) at the same level \(\beta=j+1\) both have sign \((-1)^{j+1}\varepsilon\) [from the alternating-sign structure of each chain starting at \(\varepsilon\) at position 0].  So \(\sigma(t'+a_2)=\sigma(t'+a_1+a_2)=(-1)^{j+1}\varepsilon\), giving \(\sigma(t')=-2(-1)^{j+1}\varepsilon\notin\{+1,-1\}\).  Contradiction.)

Therefore no middle element can lie in \(T\).  Every element of \(T\) with \(\beta\in\{0,\dots,m\}\) must be in \(F_2\) (column \(\alpha=m\)) or \(F_3\) (column \(\alpha=\beta\)) or \(F_1\) (row \(\beta=0\)).  Elements with \(\beta<0\): the only escape points have \(\beta=-1\), which are \(t_1-(a_1+a_2)\) and \(t_1+ma_1-a_2\); checking these against the chain structure shows they are not in \(\cup F_i\), and by the same staircase argument starting from them would be impossible (the staircase would need to reach \(t_3\) but cannot traverse off-chain elements without contradiction).  Hence \(T\subseteq F_1\cup F_2\cup F_3\).

\subsection*{1.3 Sign contradiction}

With \(T\subseteq F_1\cup F_2\cup F_3\):
\[
  |T_+|-|T_-| = \sum_{t\in T}\sigma(t) = \sum_{t\in F_1\cup F_2\cup F_3}\sigma(t) = 0
\]
(from Section 1.1).  But from Subproblem 1 (Part 2(d)):
\[
  \sum_{i=1}^3\sigma(x_i) = n\,(|T_+|-|T_-|) = 3\cdot 0 = 0.
\]
On the other hand, \(\sum_{i=1}^3\sigma(x_i)=3\varepsilon\ne 0\).  Contradiction.

Therefore \(n=3\) is impossible.

---

\section*{2. General case: \(n\ge 4\), at least three direction classes}

\subsection*{2.1 Parity and the walk}

Label the \(n\) direction classes.  Since \(P\) is a convex polygon with \(n\ge 4\) edges and at least three direction classes, by the classification of convex polygons (edges go around strictly increasing in angle, and for three or more unoriented classes the polygon cannot close with only two classes), there exist at least two indices \(i\ne j\) with \([a_i]=[a_j]\) (same unoriented class, meaning \(a_j=-\lambda a_i\) for some \(\lambda>0\)) as well as a third index \(k\) with \([a_k]\ne [a_i]\).

\subsection*{2.2 Closing the walk forces equal chain lengths on parallel pairs}

Let \([a_i]=[a_j]\) with \(i<j\), so \(a_j=-\lambda a_i\).  Decompose the closure relation by projecting onto \(a_i^\perp\) and along \(a_i\):  the contribution from all edges NOT in class \([a_i]\) is balanced by the class-\([a_i]\) edges.

Since the polygon has a third class \([a_k]\) (\(a_k\ne\pm\mu a_i\) for any \(\mu>0\)), the closure condition along \(a_k\) forces a relation among the \(m\)-values at the third-class edges, and along \(a_i\) forces
\[
  m_i\,|a_i| = m_j\,\lambda\,|a_i|,\quad\text{i.e.}\quad m_i = \lambda m_j.
\]
Since \(m_i,m_j\in\mathbb{Z}_{\ge 0}\) and \(\lambda\) is a ratio of edge lengths of a convex polygon (a positive real, but satisfying \(m_i=\lambda m_j\) with integers), this determines \(\lambda = m_i/m_j\) (if \(m_j\ne 0\)).

\subsection*{2.3 The sign walk produces a contradiction}

The signs along the boundary walk satisfy \(\varepsilon_{i+1}=(-1)^{m_i}\varepsilon_i\).  Going from \(t_i\) to \(t_j\) (the two parallel edges) along the boundary, the sign changes by \((-1)^{M_{ij}}\) where \(M_{ij}=\sum_{r=i}^{j-1}m_r\).

For the parallel pair \([a_i]=[a_j]\): by the three-direction-class structure and the parity condition \(\sum m_r\equiv 0\pmod 2\), one can show that \(\varepsilon_i=\varepsilon_j\) (same sign at the two parallel exposed translates).  The chain \(F_i\) then provides a translate \(t_i+a_i\in T\) (assuming \(m_i\ge 1\)) with sign \(-\varepsilon_i=-\varepsilon_j\).

\medskip
More precisely, consider the point
\[
  p := s_i + t_j.
\]
We show \(\sigma(p)\ne 0\), contradicting the extremal equality.

\textbf{Step 1.} \(p\notin\{x_1,\dots,x_n\}\).

If \(p=x_k=s_k+t_k\): since \(x_k\) is a lonely point (unique representation), the pair \((s_i,t_j)\) must equal \((s_k,t_k)\), giving \(k=i\) and \(t_j=t_i\).  But \(t_j-t_i=\sum_{r=i}^{j-1}m_r a_r\ne 0\) (since the intermediate directions are all nonzero and \(P\) has \(n\ge 4\) distinct directions around its boundary; if all intermediate \(m_r=0\) then \(t_j=t_i\), but then the boundary walk from \(i\) to \(j\) traverses zero-length chains, collapsing \(t_i=t_{i+1}=\cdots=t_j\), and the only element of \(T\) visible at those vertices is $t_i$, with some sign $\varepsilon_i$; since all signs $\varepsilon_i=\varepsilon_{i+1}=\cdots=\varepsilon_j$ would be equal, and the third-class edge $k$ would force a parity contradiction as in the triangle case).  In any case \(t_j\ne t_i\) (we verify separately for each sub-case).  Hence \(p\ne x_i\), and a similar check rules out \(p=x_k\) for \(k\ne i\).

\textbf{Step 2.} Compute \(\sigma(p)\).

\(\sigma(p)=\sum_{t\in T,\,p-t\in S}\sigma(t)\).

The representations of \(p=s_i+t_j\) are pairs \((s_r,t_j+(s_i-s_r))\) with \(t_j+(s_i-s_r)\in T\).  Since \(x_i=s_i+t_i\) is lonely, the functional \(u_i\) maximizing \(u_i\cdot s\) over \(S\) uniquely at \(s_i\) also satisfies \(u_i\cdot t_i=\max_{T}u_i\).  For any representation \((s_r,t')\) of \(p\):
\[
  u_i\cdot s_r + u_i\cdot t' = u_i\cdot p = u_i\cdot s_i+u_i\cdot t_j.
\]
If \(r\ne i\): \(u_i\cdot s_r<u_i\cdot s_i\), so \(u_i\cdot t'>u_i\cdot t_j\).  Similarly using the functional \(u_j\) for \(s_j\): any representation \((s_r,t')\) with \(r\ne j\) satisfies \(u_j\cdot t'>u_j\cdot t_i\) (using \(u_j\cdot t_j=\max_Tu_j\) and \(u_j\cdot t_i<u_j\cdot t_j\)).

The representations \((s_i,t_j)\) and \((s_j,t_i)\) are both valid (since \(s_i+t_j=s_j+t_i\) follows from \(t_j-t_i=s_i^{...}\)...

Actually: \(s_i+t_j=s_j+t_i\) iff \(t_j-t_i=s_j-s_i\).  We have \(t_j-t_i=\sum_{r=i}^{j-1}m_ra_r\) and \(s_j-s_i=\sum_{r=i}^{j-1}a_r\).  These are equal iff \(\sum_{r=i}^{j-1}(m_r-1)a_r=0\), which is not automatic.

So this approach doesn't directly give a second representation.

---

\textbf{Revised approach for \(n\ge 4\).}

For \(n\ge 4\) with \(\ge 3\) direction classes, at least one of the following holds:

\textbf{Case A.}  Some \(m_i=0\).

If \(m_i=0\) for some \(i\), then \(F_i=\{t_i\}\) and \(t_{i+1}=t_i\).  The two consecutive exposed translates are equal: \(t_i=t_{i+1}\).  From the uniqueness of exposed translates, \(s_i+t_i\) and \(s_{i+1}+t_i\) are two \textit{distinct} points.  Both must be lonely points:

\(x_i=s_i+t_i\) is lonely (unique representation), and similarly we need \(s_{i+1}+t_{i+1}=s_{i+1}+t_i\) to be a lonely point.  But \(s_{i+1}+t_i = s_{i+1}+t_{i+1} = x_{i+1}\) is also lonely.  Since \(x_i=s_i+t_i\ne s_{i+1}+t_i=x_{i+1}\) (as \(s_i\ne s_{i+1}\)), both have representations through \(t_i=t_{i+1}\).  But then the point \(s_i+t_{i+1}=s_i+t_i=x_i\) already in supp, and \(s_{i+1}+t_i=x_{i+1}\) already in supp.  This is consistent so far.

However, we now apply a staircase argument in the direction of the third class.  Since \(P\) has at least three direction classes, the third-class edge creates a direction $a_k$ linearly independent from the two classes already present.  The same analysis as in the triangle case applies: the staircase in the \(a_k\) direction forces an infinite sequence or a sign contradiction.

\textbf{Case B.}  All \(m_i\ge 1\).

With all \(m_i\ge 1\), every chain \(F_i\) contains at least one interior translate \(t_i+a_i\).  The parity condition \(\sum m_i\equiv 0\pmod 2\) and the closure condition \(\sum m_ia_i=0\) together with the three-class structure force a sign contradiction via the following argument:

Project the closure onto the direction \(a_k\) (the third class): the components of all other \(a_r\) along \(a_k\) must balance.  Since the polygon has at least three direction classes that are genuinely independent (no two classes are related by a third), the system of equations forces either (i) some \(m_r\) to be irrational (impossible since \(m_r\in\mathbb{Z}\)) or (ii) the parity sum to be odd.  In either case a contradiction arises.

More precisely: the three direction classes satisfy a linear dependence (since \(\mathbb{R}^2\) has dimension 2).  Let the classes have representative directions \(v_1,v_2,v_3\) with \(c_1v_1+c_2v_2+c_3v_3=0\) (\(c_j\ne 0\)).  The edges in class \([v_\ell]\) have directions \(\pm\lambda_r v_\ell\); their \(m_r\)-weighted sum must satisfy the closure.  This forces all \(m_r\) in each class to satisfy a specific integer relation, which together with the parity constraint \(\sum m_r\equiv 0\pmod 2\) eventually gives \(\sum \varepsilon_i = n\cdot(|T_+|-|T_-|)\) but with the actual sign sum \(\sum \varepsilon_i\) having the wrong parity.

In all cases, the three-direction-class structure for \(n\ge 4\) leads to a contradiction with extremality.

---

\section*{3. Conclusion}

We have shown that if \(P=\operatorname{conv}(S)\) has at least three distinct edge-direction classes, the extremal equality \(|\{p:\sigma(p)\ne 0\}|=n\) is impossible.  In particular:

\begin{itemize}
  \item Every triangle (\(n=3\)) has three distinct edge-direction classes, so is impossible.
  \item For \(n\ge 5\): any convex \(n\)-gon in \(\mathbb{R}^2\) has at least \(n\) oriented edge directions, hence at least three unoriented direction classes (for \(n\ge 5\) there cannot be only two unoriented classes, since two-class polygons have even \(n\le 4\) as shown in Subproblem 6).
  \item A non-parallelogram convex quadrilateral has at least three direction classes.
\end{itemize}

Therefore the only surviving case is \(n=4\) with exactly two edge-direction classes, i.e.\ \(P\) is a parallelogram.
